\documentclass[12pt, a4paper]{article}

\usepackage[utf8]{inputenc}
\usepackage[T1]{fontenc}
\usepackage{helvet}
\renewcommand{\familydefault}{\sfdefault}
\usepackage{geometry}
\setlength{\headheight}{15pt}
\geometry{top=2.5cm, bottom=2.5cm, left=2.5cm, right=2.5cm}

\usepackage{array}
\usepackage{booktabs}
\usepackage{longtable}
\usepackage{xcolor}
\usepackage{colortbl}
\usepackage{tabularx}
\usepackage{titlesec}
\usepackage{fancyhdr}
\usepackage{enumitem}
\usepackage{hyperref}
\usepackage{parskip}
\usepackage{multirow}
\usepackage{mdframed}
\usepackage{tcolorbox}
\tcbuselibrary{skins, breakable}

% ── Colores ───────────────────────────────────────────────────────────────────
\definecolor{azulOscuro}{RGB}{30, 58, 95}
\definecolor{azulMedio}{RGB}{52, 100, 163}
\definecolor{azulClaro}{RGB}{209, 226, 245}
\definecolor{grisFondo}{RGB}{248, 249, 251}
\definecolor{grisLinea}{RGB}{220, 224, 230}

% ── Encabezado/pie ────────────────────────────────────────────────────────────
\pagestyle{fancy}
\fancyhf{}
\fancyhead[L]{\small\color{azulOscuro}\textbf{Sistema de Vigilancia Docente}}
\fancyhead[R]{\small\color{azulOscuro}Casos de Uso -- Primera Entrega}
\fancyfoot[C]{\small\color{gray}\thepage}
\renewcommand{\headrulewidth}{0.6pt}

% ── Títulos ───────────────────────────────────────────────────────────────────
\titleformat{\section}{\Large\bfseries\color{azulOscuro}}{\thesection.}{0.6em}{}
  [\vspace{-2pt}\color{azulMedio}\rule{\linewidth}{1.5pt}\vspace{2pt}]

\titleformat{\subsection}{\normalsize\bfseries\color{azulMedio}}{\thesubsection}{0.5em}{}

% ── Caja de caso de uso ───────────────────────────────────────────────────────
\tcbset{
  casousostyle/.style={
    enhanced, breakable,
    colback=grisFondo,
    colframe=azulOscuro,
    boxrule=1.2pt,
    arc=4pt,
    left=10pt, right=10pt, top=8pt, bottom=8pt,
    fonttitle=\bfseries\large\color{white},
    attach boxed title to top left={yshift=-2mm, xshift=6mm},
    boxed title style={colback=azulOscuro, arc=3pt, boxrule=0pt},
  }
}

% Comando: \cu{ID -- Nombre}{Actor}{ contenido tabular }
\newcommand{\cu}[3]{%
  \begin{tcolorbox}[casousostyle, title={#1},
    after title={\hfill\small\normalfont\color{white!80!azulOscuro}#2}]
  \renewcommand{\arraystretch}{1.35}
  \begin{tabularx}{\linewidth}{>{\bfseries\color{azulOscuro}\small}p{3.6cm} X}
  #3
  \end{tabularx}
  \end{tcolorbox}
  \vspace{6pt}
}

\newcommand{\fcu}[2]{\small\textbf{\color{azulOscuro}#1} & \small #2 \\}

% ── Pasos numerados ───────────────────────────────────────────────────────────
\newlist{pasos}{enumerate}{1}
\setlist[pasos]{label=\textbf{\arabic*.}, leftmargin=*, itemsep=2pt, topsep=2pt, parsep=0pt}

% ── Hyperref ──────────────────────────────────────────────────────────────────
\hypersetup{colorlinks=true, linkcolor=azulMedio, urlcolor=azulMedio,
  pdftitle={Casos de Uso -- Sistema de Vigilancia Docente}}

% ══════════════════════════════════════════════════════════════════════════════
\begin{document}
\thispagestyle{empty}

% ── Portada ───────────────────────────────────────────────────────────────────
\begin{center}
  \vspace*{1.5cm}
  {\color{azulOscuro}\rule{\linewidth}{2.5pt}}
  \vspace{0.5cm}

  {\Huge\bfseries\color{azulOscuro} Sistema Hibrido para la\\[6pt] Vigilancia Docente}

  \vspace{0.4cm}
  {\color{azulMedio}\rule{\linewidth}{1pt}}
  \vspace{0.6cm}

  {\Large\color{azulMedio} Descripcion Detallada de Casos de Uso}

  \vspace{1.2cm}

  \begin{tcolorbox}[colback=azulClaro, colframe=azulMedio, arc=6pt,
    boxrule=1pt, left=16pt, right=16pt, top=12pt, bottom=12pt]
    \centering
    {\large\bfseries\color{azulOscuro} Proyecto de Desarrollo Web 2026}\\[6pt]
    {\normalsize\color{azulOscuro} Ing.\ Juan Camilo Gonzalez Vargas}\\[6pt]
    {\normalsize\color{azulOscuro} Primera Entrega --- 10\%}
  \end{tcolorbox}

  \vspace{1.5cm}

  {\normalsize\color{gray}
    \begin{tabular}{ll}
      \textbf{Version:} & 1.0 \\[2pt]
      \textbf{Fecha:}   & Marzo 2026 \\[2pt]
      \textbf{Estado:}  & Primera Entrega \\
    \end{tabular}
  }

  \vfill
  {\color{azulOscuro}\rule{\linewidth}{2.5pt}}
\end{center}

\newpage
\tableofcontents
\newpage

% ══════════════════════════════════════════════════════════════════════════════
\section{Introduccion}

Este documento presenta la descripcion detallada de los casos de uso del sistema
hibrido de vigilancia docente durante los periodos de recreo y almuerzo. El sistema
busca garantizar la presencia puntual del docente, promover recorridos activos y
generar analitica accionable para la prevencion de incidentes.

\subsection{Actores del sistema}

\renewcommand{\arraystretch}{1.4}
\begin{tabularx}{\linewidth}{>{\bfseries\color{azulOscuro}}p{3.8cm} X}
  \toprule
  \rowcolor{azulClaro}
  \color{azulOscuro}Actor & \color{azulOscuro}Descripcion \\
  \midrule
  Docente en turno
    & Recibe notificaciones, realiza check-in, registra recorridos e incidentes,
      califica la limpieza y puede solicitar reasignacion por impedimento. \\[2pt]
  \rowcolor{grisFondo}
  Coordinador de nivel
    & Visualiza el tablero en tiempo real, recibe alertas por zonas sin cobertura,
      gestiona reasignaciones y consulta reportes y analitica. \\[2pt]
  Administrador
    & Define zonas y checkpoints, carga turnos y asignaciones, y configura reglas
      operativas (umbrales, escalamiento, gamificacion). \\
  \bottomrule
\end{tabularx}

\subsection{Convenciones}

\begin{itemize}[leftmargin=*, itemsep=2pt]
  \item \textbf{CU-XX}: identificador unico del caso de uso.
  \item \textbf{Flujo principal}: secuencia de pasos del escenario exitoso.
  \item \textbf{Flujos alternativos}: variaciones o excepciones del flujo principal.
  \item \textbf{Precondicion}: estado del sistema antes de ejecutar el caso de uso.
  \item \textbf{Postcondicion}: estado del sistema tras la ejecucion exitosa.
\end{itemize}

\newpage

% ══════════════════════════════════════════════════════════════════════════════
\section{Modulo de Turnos y Notificaciones}

\cu{CU-01 --- Asignar turno de vigilancia}{Administrador}{
  \fcu{Actor principal}{Administrador del sistema}
  \fcu{Actor secundario}{Docente (notificado)}
  \fcu{Precondicion}{El administrador ha iniciado sesion. Existe al menos un docente activo y una zona activa.}
  \fcu{Postcondicion}{El turno queda registrado con estado \textbf{PENDIENTE} y el docente recibe una notificacion de tipo RECORDATORIO.}
  \fcu{Flujo principal}{
    \begin{pasos}
      \item El administrador accede al modulo de Turnos.
      \item Selecciona ``Asignar turno''.
      \item Escoge el docente, la zona, la franja horaria (RECREO / ALMUERZO), la fecha y la hora de inicio y fin.
      \item Confirma el formulario.
      \item El sistema crea el turno con estado PENDIENTE.
      \item El sistema genera una notificacion de tipo RECORDATORIO para el docente asignado.
    \end{pasos}
  }
  \fcu{Flujos alternativos}{
    \textbf{FA1 --- Docente no disponible:} Si el docente ya tiene un turno activo en la misma franja, el sistema advierte y sugiere otro docente disponible. \\[4pt]
    \textbf{FA2 --- Datos incompletos:} Si faltan campos obligatorios, el sistema muestra un error y no guarda el registro.
  }
  \fcu{Reglas de negocio}{Un docente no puede tener dos turnos activos en la misma franja horaria y zona.}
}

\cu{CU-02 --- Visualizar calendario de turnos}{Docente / Coordinador / Admin}{
  \fcu{Actor principal}{Docente en turno}
  \fcu{Precondicion}{El usuario ha iniciado sesion.}
  \fcu{Postcondicion}{El usuario visualiza los turnos asignados con su estado actual.}
  \fcu{Flujo principal}{
    \begin{pasos}
      \item El usuario accede al modulo de Turnos.
      \item El sistema lista los turnos con: docente asignado, zona, franja, fecha/hora y estado (PENDIENTE / EN\_CURSO / CERRADO).
      \item El usuario puede filtrar por fecha o zona.
      \item El coordinador visualiza todos los turnos del nivel.
    \end{pasos}
  }
  \fcu{Flujos alternativos}{
    \textbf{FA1 --- Sin turnos:} El sistema muestra un mensaje indicando que no hay turnos para el filtro seleccionado.
  }
  \fcu{Reglas de negocio}{El docente solo ve sus propios turnos. El coordinador y el administrador ven todos.}
}

\cu{CU-03 --- Recibir recordatorio de turno}{Sistema / Docente}{
  \fcu{Actor principal}{Sistema (automatico)}
  \fcu{Actor secundario}{Docente en turno}
  \fcu{Precondicion}{Existe un turno en estado PENDIENTE proximo a iniciar.}
  \fcu{Postcondicion}{El docente recibe notificaciones escalonadas con la informacion del turno.}
  \fcu{Flujo principal}{
    \begin{pasos}
      \item El sistema detecta un turno que iniciara en 10 minutos.
      \item Envia una notificacion de tipo RECORDATORIO al docente con: zona asignada, hora exacta y boton ``Abrir turno''.
      \item A los 5 minutos, el sistema envia un segundo recordatorio de urgencia.
      \item El docente visualiza la notificacion en su bandeja.
    \end{pasos}
  }
  \fcu{Flujos alternativos}{
    \textbf{FA1 --- Turno ya iniciado:} Si el docente abre el recordatorio despues del inicio, el sistema muestra directamente el boton de check-in.
  }
  \fcu{Reglas de negocio}{Las notificaciones se generan automaticamente a los 10 y 5 minutos previos al inicio del turno.}
}

\newpage

% ══════════════════════════════════════════════════════════════════════════════
\section{Modulo de Verificacion de Presencia}

\cu{CU-04 --- Realizar check-in de turno}{Docente}{
  \fcu{Actor principal}{Docente en turno}
  \fcu{Precondicion}{Existe un turno en estado PENDIENTE asignado al docente. El docente esta fisicamente en la zona.}
  \fcu{Postcondicion}{El turno pasa a estado \textbf{EN\_CURSO}. El tablero del coordinador marca la zona como cubierta (verde).}
  \fcu{Flujo principal}{
    \begin{pasos}
      \item El docente accede al modulo de Check-ins.
      \item Selecciona el turno activo pendiente.
      \item Escanea el codigo QR fijo del punto fisico de la zona.
      \item El sistema valida el QR y registra el check-in con la marca de tiempo.
      \item El turno cambia a estado EN\_CURSO.
      \item El tablero actualiza la zona a estado verde (cubierta).
    \end{pasos}
  }
  \fcu{Flujos alternativos}{
    \textbf{FA1 --- Usar PIN:} Si el QR no esta disponible, el docente ingresa el PIN rotativo. El sistema valida y registra el check-in.\\[4pt]
    \textbf{FA2 --- QR invalido:} El sistema rechaza el check-in y solicita reintento.\\[4pt]
    \textbf{FA3 --- Turno ya iniciado:} El sistema notifica que el turno ya se encuentra en curso.
  }
  \fcu{Reglas de negocio}{Solo se puede hacer check-in dentro de la ventana de tiempo del turno. QR es el metodo principal; PIN y NFC son alternativas.}
}

\cu{CU-05 --- Alertar por ausencia de cobertura}{Sistema / Coordinador}{
  \fcu{Actor principal}{Sistema (automatico)}
  \fcu{Actor secundario}{Coordinador de nivel}
  \fcu{Precondicion}{Un turno supera el umbral de tiempo (2 min) sin check-in.}
  \fcu{Postcondicion}{El coordinador recibe una alerta. La zona figura en rojo en el tablero.}
  \fcu{Flujo principal}{
    \begin{pasos}
      \item El sistema detecta un turno sin check-in despues de 2 minutos del inicio.
      \item Cambia el estado visual de la zona a rojo (sin cobertura) en el tablero.
      \item Genera una notificacion de tipo ALERTA al coordinador con: zona, docente asignado y tiempo transcurrido.
      \item El coordinador visualiza la alerta con el boton de accion rapida ``Reasignar''.
    \end{pasos}
  }
  \fcu{Flujos alternativos}{
    \textbf{FA1 --- Docente llega tarde:} Si el docente hace check-in despues del umbral, el sistema registra el retraso, actualiza la zona a verde y mantiene la alerta como historico.
  }
  \fcu{Reglas de negocio}{El umbral de ausencia es configurable por el administrador. La reasignacion debe completarse en maximo 2 interacciones.}
}

\newpage

% ══════════════════════════════════════════════════════════════════════════════
\section{Modulo de Vigilancia Activa}

\cu{CU-06 --- Registrar recorrido activo}{Docente}{
  \fcu{Actor principal}{Docente en turno}
  \fcu{Precondicion}{El turno esta en estado EN\_CURSO. Han transcurrido X minutos sin evidencia de recorrido.}
  \fcu{Postcondicion}{Se registra un CheckIn de tipo recorrido (\texttt{esRecorrido=true}). La metrica de recorridos del docente se actualiza.}
  \fcu{Flujo principal}{
    \begin{pasos}
      \item El sistema detecta que el docente lleva mas de X minutos sin evidencia de movimiento.
      \item El sistema muestra el recordatorio ``Confirmar recorrido'' en pantalla.
      \item El docente se desplaza hasta un checkpoint de la zona.
      \item Escanea el QR del checkpoint.
      \item El sistema registra el check-in con \texttt{esRecorrido=true} y actualiza la evidencia de movilidad.
    \end{pasos}
  }
  \fcu{Flujos alternativos}{
    \textbf{FA1 --- Confirmacion manual:} Si no hay checkpoint disponible, el docente confirma el recorrido manualmente. El sistema lo registra distinguiendolo del escaneo fisico.\\[4pt]
    \textbf{FA2 --- Sin respuesta:} Si el docente no responde al recordatorio, el coordinador recibe una alerta de vigilancia pasiva.
  }
  \fcu{Reglas de negocio}{La confirmacion manual no reemplaza sistematicamente el escaneo fisico. Cada zona debe tener al menos un checkpoint.}
}

\newpage

% ══════════════════════════════════════════════════════════════════════════════
\section{Modulo de Registro de Incidentes}

\cu{CU-07 --- Registrar incidente}{Docente}{
  \fcu{Actor principal}{Docente en turno}
  \fcu{Precondicion}{El docente tiene un turno en estado EN\_CURSO.}
  \fcu{Postcondicion}{El incidente queda registrado con tipo, severidad, zona y marca de tiempo. Si la severidad es S3, el coordinador recibe una alerta inmediata.}
  \fcu{Flujo principal}{
    \begin{pasos}
      \item El docente accede al boton ``Registrar situacion'' desde la pantalla de turno.
      \item Selecciona el tipo de incidente mediante lista desplegable:
        \begin{itemize}[itemsep=1pt, topsep=2pt]
          \item Seguridad fisica (caida, golpe, accidente leve)
          \item Convivencia (pelea, agresion verbal, intimidacion)
          \item Uso del espacio (mobiliario, infraestructura, juego peligroso)
          \item Observacion social
        \end{itemize}
      \item Selecciona la severidad: S1 (leve), S2 (requiere seguimiento) o S3 (atencion inmediata).
      \item Agrega una descripcion breve (opcional).
      \item Si el tipo es ``Observacion social'', ingresa el curso del estudiante.
      \item Confirma el registro.
      \item El sistema guarda el incidente con la zona del turno y la marca de tiempo actual.
    \end{pasos}
  }
  \fcu{Flujos alternativos}{
    \textbf{FA1 --- Severidad S3:} El sistema genera inmediatamente una notificacion de tipo INCIDENTE al coordinador.\\[4pt]
    \textbf{FA2 --- Sin turno activo:} Si el docente no tiene turno en curso, puede registrar el incidente asociandolo manualmente a una zona.
  }
  \fcu{Reglas de negocio}{No se registra el nombre del estudiante (proteccion de datos). Solo se admite el curso en el tipo SOCIAL. El registro debe completarse en maximo 3 pasos.}
}

\cu{CU-08 --- Consultar historial de incidentes}{Coordinador / Admin}{
  \fcu{Actor principal}{Coordinador de nivel}
  \fcu{Precondicion}{El coordinador ha iniciado sesion.}
  \fcu{Postcondicion}{El coordinador visualiza el listado de incidentes filtrado segun sus criterios.}
  \fcu{Flujo principal}{
    \begin{pasos}
      \item El coordinador accede al modulo de Incidentes.
      \item Aplica filtros: tipo, severidad, zona o franja horaria.
      \item El sistema muestra el listado con tipo, severidad, zona, docente que reporto y timestamp.
      \item El coordinador puede seleccionar un incidente para ver su detalle completo.
    \end{pasos}
  }
  \fcu{Flujos alternativos}{
    \textbf{FA1 --- Sin resultados:} El sistema informa que no existen incidentes para el filtro aplicado.
  }
  \fcu{Reglas de negocio}{El coordinador solo ve incidentes de su nivel. El administrador tiene visibilidad global.}
}

\newpage

% ══════════════════════════════════════════════════════════════════════════════
\section{Modulo de Reasignacion}

\cu{CU-09 --- Solicitar reasignacion por impedimento}{Docente}{
  \fcu{Actor principal}{Docente en turno}
  \fcu{Precondicion}{El docente tiene un turno en estado PENDIENTE o EN\_CURSO y no puede cumplirlo.}
  \fcu{Postcondicion}{Se crea una Reasignacion en estado PROPUESTA. El docente candidato recibe una notificacion.}
  \fcu{Flujo principal}{
    \begin{pasos}
      \item El docente presiona el boton ``No puedo llegar al turno''.
      \item Ingresa el motivo del impedimento.
      \item El sistema busca candidatos: sin turno en esa franja, con menor carga y cercanos a la zona.
      \item El sistema propone el candidato mas adecuado automaticamente.
      \item El candidato recibe una notificacion de tipo REASIGNACION con ventana de aceptacion de 30 a 60 segundos.
      \item El candidato acepta o rechaza.
      \item Si acepta, la reasignacion pasa a estado ACEPTADA y el turno queda asignado al nuevo docente.
      \item El coordinador recibe notificacion del cambio.
    \end{pasos}
  }
  \fcu{Flujos alternativos}{
    \textbf{FA1 --- Candidato rechaza:} El sistema propone al siguiente candidato disponible y reinicia la ventana de aceptacion.\\[4pt]
    \textbf{FA2 --- Sin candidatos:} El coordinador recibe una alerta critica para gestionar manualmente.\\[4pt]
    \textbf{FA3 --- Ventana expirada:} El sistema propone automaticamente al siguiente candidato.
  }
  \fcu{Reglas de negocio}{Trazabilidad completa: se registra quien reasigno, motivo, hora de propuesta y hora de aceptacion/rechazo.}
}

\cu{CU-10 --- Gestionar reasignacion}{Coordinador}{
  \fcu{Actor principal}{Coordinador de nivel}
  \fcu{Precondicion}{Existe una reasignacion en estado PROPUESTA o sin candidatos disponibles.}
  \fcu{Postcondicion}{El coordinador resuelve la cobertura de la zona afectada.}
  \fcu{Flujo principal}{
    \begin{pasos}
      \item El coordinador visualiza la alerta de reasignacion en el tablero.
      \item Accede al modulo de Reasignaciones.
      \item Revisa el turno afectado, el docente original y los candidatos propuestos.
      \item Asigna manualmente un docente de reemplazo si la propuesta automatica fallo.
      \item Confirma la reasignacion.
      \item El sistema actualiza el turno y notifica al docente de reemplazo.
    \end{pasos}
  }
  \fcu{Flujos alternativos}{
    \textbf{FA1 --- Reasignacion ya resuelta:} El sistema informa que la reasignacion fue aceptada automaticamente.
  }
  \fcu{Reglas de negocio}{El coordinador puede sobrescribir la propuesta automatica del sistema en cualquier momento.}
}

\newpage

% ══════════════════════════════════════════════════════════════════════════════
\section{Modulo de Registro de Limpieza}

\cu{CU-11 --- Registrar limpieza al cierre del turno}{Docente}{
  \fcu{Actor principal}{Docente en turno}
  \fcu{Precondicion}{El turno esta en estado EN\_CURSO. El docente se dispone a cerrarlo.}
  \fcu{Postcondicion}{Se registra el estado de limpieza de la zona. El turno pasa a estado \textbf{CERRADO}.}
  \fcu{Flujo principal}{
    \begin{pasos}
      \item El docente accede al modulo de cierre de turno.
      \item El sistema solicita la escala de limpieza obligatoria:
        \begin{itemize}[itemsep=1pt, topsep=2pt]
          \item 1 --- Limpio
          \item 2 --- Algo de basura
          \item 3 --- Mucha basura
          \item 4 --- Critico
        \end{itemize}
      \item El docente selecciona la escala correspondiente.
      \item Agrega una observacion opcional.
      \item Confirma el cierre.
      \item El sistema guarda el RegistroLimpieza y cambia el turno a estado CERRADO.
    \end{pasos}
  }
  \fcu{Flujos alternativos}{
    \textbf{FA1 --- Escala 3 o 4:} El sistema genera una notificacion al coordinador indicando el estado critico de limpieza de la zona.
  }
  \fcu{Reglas de negocio}{La escala de limpieza es \textbf{obligatoria} para cerrar el turno. Solo puede existir un registro de limpieza por turno (relacion 1:1).}
}

\newpage

% ══════════════════════════════════════════════════════════════════════════════
\section{Modulo de Analitica}

\cu{CU-12 --- Consultar mapa de calor de incidentes}{Coordinador / Admin}{
  \fcu{Actor principal}{Coordinador de nivel}
  \fcu{Precondicion}{Existen incidentes registrados en el sistema.}
  \fcu{Postcondicion}{El coordinador visualiza la distribucion de incidentes por zona, franja y tipo.}
  \fcu{Flujo principal}{
    \begin{pasos}
      \item El coordinador accede al modulo de Analitica --- Mapa de Calor.
      \item Selecciona el rango de semanas a analizar.
      \item El sistema muestra una tabla con el porcentaje de incidentes por zona, desglosado por tipo (Fisico, Convivencia, Espacio, Social) y franja (Recreo, Almuerzo).
      \item Las celdas se colorean segun la intensidad: verde (0), amarillo (1), naranja (2--3), rojo (4+).
      \item El coordinador identifica las zonas criticas y los patrones temporales.
    \end{pasos}
  }
  \fcu{Flujos alternativos}{
    \textbf{FA1 --- Sin datos:} El sistema muestra la tabla en blanco con un mensaje informativo.
  }
  \fcu{Reglas de negocio}{Los datos del mapa de calor se calculan en tiempo real a partir de los incidentes registrados.}
}

\cu{CU-13 --- Consultar tablero en vivo}{Coordinador}{
  \fcu{Actor principal}{Coordinador de nivel}
  \fcu{Precondicion}{El coordinador ha iniciado sesion. Existen turnos en la franja horaria activa.}
  \fcu{Postcondicion}{El coordinador visualiza el estado de cobertura de todas las zonas en tiempo real.}
  \fcu{Flujo principal}{
    \begin{pasos}
      \item El coordinador accede al Dashboard.
      \item El sistema muestra el tablero con el estado de cada zona:
        \begin{itemize}[itemsep=1pt, topsep=2pt]
          \item \textbf{Verde}: zona cubierta (check-in realizado)
          \item \textbf{Amarillo}: turno proximo a iniciar (en ventana de espera)
          \item \textbf{Rojo}: zona sin cobertura (sin check-in tras el umbral)
        \end{itemize}
      \item El coordinador puede filtrar por franja horaria o zona.
      \item Por cada zona ve: docente asignado, hora del check-in, estado de recorrido y acceso rapido a ``Reasignar''.
    \end{pasos}
  }
  \fcu{Flujos alternativos}{
    \textbf{FA1 --- Zona roja:} El coordinador presiona ``Reasignar'' directamente desde el tablero para activar el flujo CU-10 en maximo 2 toques.
  }
  \fcu{Reglas de negocio}{El tablero debe reflejar cambios en menos de 3 segundos. Las acciones criticas se resuelven en maximo 2 interacciones.}
}

\cu{CU-14 --- Ver metricas de gamificacion}{Docente / Coordinador}{
  \fcu{Actor principal}{Docente en turno}
  \fcu{Precondicion}{El docente tiene al menos un trimestre de datos registrados.}
  \fcu{Postcondicion}{El docente visualiza sus indicadores positivos y su posicion en el ranking trimestral.}
  \fcu{Flujo principal}{
    \begin{pasos}
      \item El docente accede al modulo de Metricas.
      \item El sistema muestra sus indicadores del trimestre: puntualidad (\%), total de recorridos, calidad del registro (\%), contribucion preventiva (\%) y puntaje total.
      \item Si el docente tiene reconocimiento trimestral, el sistema lo destaca con una insignia.
      \item El coordinador puede ver el ranking completo de todos los docentes del nivel.
    \end{pasos}
  }
  \fcu{Flujos alternativos}{
    \textbf{FA1 --- Sin datos:} El sistema muestra cero en todos los indicadores con un mensaje motivacional.
  }
  \fcu{Reglas de negocio}{El enfoque es no punitivo: solo se muestran metricas positivas. Los reconocimientos se otorgan trimestralmente.}
}

\newpage

% ══════════════════════════════════════════════════════════════════════════════
\section{Modulo de Configuracion}

\cu{CU-15 --- Gestionar zonas y checkpoints}{Administrador}{
  \fcu{Actor principal}{Administrador del sistema}
  \fcu{Precondicion}{El administrador ha iniciado sesion.}
  \fcu{Postcondicion}{La zona o checkpoint queda creado, modificado o desactivado.}
  \fcu{Flujo principal}{
    \begin{pasos}
      \item El administrador accede al modulo de Zonas.
      \item Selecciona ``Nueva zona'' o edita una existente.
      \item Ingresa nombre, descripcion, capacidad, codigo QR y PIN rotativo.
      \item Guarda la zona.
      \item Para agregar checkpoints, accede al modulo de Checkpoints, selecciona la zona y crea los puntos de recorrido con su codigo QR individual.
    \end{pasos}
  }
  \fcu{Flujos alternativos}{
    \textbf{FA1 --- QR duplicado:} El sistema rechaza el registro y solicita un codigo unico.\\[4pt]
    \textbf{FA2 --- Desactivar zona con turnos activos:} El sistema advierte antes de desactivar.
  }
  \fcu{Reglas de negocio}{Una zona desactivada no puede recibir nuevos turnos pero conserva el historial de incidentes y metricas.}
}

\cu{CU-16 --- Gestionar usuarios}{Administrador}{
  \fcu{Actor principal}{Administrador del sistema}
  \fcu{Precondicion}{El administrador ha iniciado sesion.}
  \fcu{Postcondicion}{El usuario queda creado, modificado o desactivado con el rol correcto.}
  \fcu{Flujo principal}{
    \begin{pasos}
      \item El administrador accede al modulo de Usuarios.
      \item Selecciona ``Nuevo usuario''.
      \item Ingresa nombre, correo institucional, contrasena y rol (DOCENTE, COORDINADOR o ADMIN).
      \item Guarda el registro.
      \item El nuevo usuario puede iniciar sesion con sus credenciales.
    \end{pasos}
  }
  \fcu{Flujos alternativos}{
    \textbf{FA1 --- Email duplicado:} El sistema rechaza el registro e informa que el correo ya existe.\\[4pt]
    \textbf{FA2 --- Desactivar usuario:} Los turnos futuros quedan en PENDIENTE y se notifica al coordinador para reasignarlos.
  }
  \fcu{Reglas de negocio}{El correo institucional es unico e irrepetible. Un usuario inactivo no recibe notificaciones ni puede iniciar sesion.}
}

\newpage

% ══════════════════════════════════════════════════════════════════════════════
\section{Resumen de Casos de Uso}

\renewcommand{\arraystretch}{1.4}
\begin{longtable}{>{\bfseries\small\color{azulOscuro}}p{1.6cm} >{\small}p{6cm} >{\small}p{3cm} >{\small}p{3cm}}
  \toprule
  \rowcolor{azulOscuro}
  \color{white}ID &
  \color{white}Nombre &
  \color{white}Actor principal &
  \color{white}Modulo \\
  \midrule
  \endfirsthead
  \toprule
  \rowcolor{azulOscuro}
  \color{white}ID &
  \color{white}Nombre &
  \color{white}Actor principal &
  \color{white}Modulo \\
  \midrule
  \endhead
  \bottomrule
  \endfoot

  CU-01 & Asignar turno de vigilancia & Administrador & Turnos \\
  \rowcolor{grisFondo}
  CU-02 & Visualizar calendario de turnos & Docente / Coord. & Turnos \\
  CU-03 & Recibir recordatorio de turno & Sistema & Notificaciones \\
  \rowcolor{grisFondo}
  CU-04 & Realizar check-in de turno & Docente & Verificacion \\
  CU-05 & Alertar por ausencia de cobertura & Sistema / Coord. & Verificacion \\
  \rowcolor{grisFondo}
  CU-06 & Registrar recorrido activo & Docente & Vigilancia activa \\
  CU-07 & Registrar incidente & Docente & Incidentes \\
  \rowcolor{grisFondo}
  CU-08 & Consultar historial de incidentes & Coordinador & Incidentes \\
  CU-09 & Solicitar reasignacion por impedimento & Docente & Reasignacion \\
  \rowcolor{grisFondo}
  CU-10 & Gestionar reasignacion & Coordinador & Reasignacion \\
  CU-11 & Registrar limpieza al cierre & Docente & Limpieza \\
  \rowcolor{grisFondo}
  CU-12 & Consultar mapa de calor & Coordinador & Analitica \\
  CU-13 & Consultar tablero en vivo & Coordinador & Analitica \\
  \rowcolor{grisFondo}
  CU-14 & Ver metricas de gamificacion & Docente / Coord. & Analitica \\
  CU-15 & Gestionar zonas y checkpoints & Administrador & Configuracion \\
  \rowcolor{grisFondo}
  CU-16 & Gestionar usuarios & Administrador & Configuracion \\
\end{longtable}

\end{document}
